\section{Configuration}
\label{sec:configuration}

This section covers the runtime configuration surfaces: the \raft{} cluster
Compose file and DNS discovery, the \ipto{} backend environment variables, the
PROV-O SDL schema, attribute mapping, and the \raft{}-side control parameters.

\subsection{Raft Cluster Configuration}

The three-node cluster is configured entirely through
\filepath{docker-compose.cluster.yml} and the CoreDNS zone. Each node is given a
static address on the \texttt{10.77.0.0/24} bridge network and a command line of
the form:

\begin{lstlisting}[numbers=none]
--srv _vannak._tcp.vannak.local 0.0.0.0 10081 vannak-1 /data
\end{lstlisting}

The arguments are the SRV service name to resolve, the bind host, the listen
port, the node id, and the data directory. The Compose file points each node's
DNS resolver at the CoreDNS container (\texttt{10.77.0.10}).

\subsection{Static IP Addresses}

Static addressing makes the DNS zone deterministic and avoids relying on Docker
DNS for the SRV records.

\begin{table}[htbp]
\centering
\begin{tabular}{ll}
\hline
Host & Address \\
\hline
\texttt{ns.vannak.local} & \texttt{10.77.0.10} \\
\texttt{vannak-1.vannak.local} & \texttt{10.77.0.11} \\
\texttt{vannak-2.vannak.local} & \texttt{10.77.0.12} \\
\texttt{vannak-3.vannak.local} & \texttt{10.77.0.13} \\
\hline
\end{tabular}
\caption{Static address assignments in the cluster zone.}
\end{table}

\subsection{CoreDNS SRV Records}

CoreDNS serves the \texttt{vannak.local} zone from \filepath{dns/vannak.local.db}
and forwards everything else upstream. The Corefile declares both zones:

\begin{lstlisting}[numbers=none]
vannak.local:53 {
    file /vannak.local.db
    log
    errors
}

.:53 {
    forward . 1.1.1.1 8.8.8.8
    log
    errors
}
\end{lstlisting}

The zone file publishes one A-record per node and three SRV records under
\texttt{\_vannak.\_tcp}:

\begin{lstlisting}[numbers=none]
_vannak._tcp IN SRV 10 10 10081 vannak-1.vannak.local.
_vannak._tcp IN SRV 10 10 10081 vannak-2.vannak.local.
_vannak._tcp IN SRV 10 10 10081 vannak-3.vannak.local.
\end{lstlisting}

\subsection{DNS Auto-Discovery}

At startup, \texttt{vannak-node} resolves its SRV service by shelling out to
\shellcmd{dig +short SRV _vannak._tcp.vannak.local}. Each returned record yields
a peer specification of the form \texttt{peer-id@hostname:port}, where the peer
id is the first DNS label of the hostname. If no peers resolve, the node logs a
warning and starts as a single node.

\subsection{Port 10081}

All cluster nodes listen on TCP port 10081 for the \graft{} transport. The same
port is used for the \texttt{probe} subcommand. The Dockerfile declares
\shellcmd{EXPOSE 10081}.

\subsection{Ipto Backend Environment Variables}

The \texttt{ipto-writer} feature connects to PostgreSQL through
\texttt{ipto\_rust}'s \texttt{PostgresBackend}, which reads \texttt{IPTO\_PG\_*}
environment variables with sensible local defaults.

\begin{table}[htbp]
\centering
\begin{tabular}{ll}
\hline
Variable & Default \\
\hline
\texttt{IPTO\_PG\_HOST} & \texttt{localhost} \\
\texttt{IPTO\_PG\_PORT} & \texttt{5432} \\
\texttt{IPTO\_PG\_USER} & \texttt{repo} \\
\texttt{IPTO\_PG\_PASSWORD} & \texttt{repo} \\
\texttt{IPTO\_PG\_DATABASE} & \texttt{repo} \\
\hline
\end{tabular}
\caption{\ipto{} PostgreSQL backend environment variables.}
\end{table}

These defaults match the credentials set by \filepath{docker-compose.test.yml},
so the integration test connects with no extra configuration.

\subsection{The PROV-O SDL Schema}

\texttt{IptoRepoWriter} ships a built-in SDL schema, the constant
\texttt{PROV\_O\_SDL}, declaring the provenance attribute registry and a
\texttt{ProvEntity} unit template. It registers a datatype registry (string,
time, integer, long, double, boolean), the core PROV-O properties, RDFS label
and comment, and \vannak{}-specific attributes such as
\texttt{vannak:dataIndividualId} and \texttt{vannak:processInstanceId}. The
schema is persisted by calling \texttt{configure\_sdl()} once at startup.

\subsection{Attribute Mapping}

An \texttt{IptoMapping} declares a version string and a set of mappings from
\texttt{MetadataFieldName} to \texttt{IptoAttributeName}. Only mapped fields are
written; unmapped passive or active fields are ignored when building the
payload.

\begin{lstlisting}
let mapping = IptoMapping::new("v1")
    .map_field("metadata.source_system", "attr:provenance.source_system")
    .map_field("mask.customer.email", "attr:provenance.masked_field");
\end{lstlisting}

\subsection{Raft Parameters}

The \texttt{vannak-node} binary constructs a \texttt{RaftNode} with the
discovered membership and a fixed set of tuning arguments (an election tick
base, a snapshot/log threshold, and a maximum batch size). These are passed
positionally to \texttt{RaftNode::new}; adjust them in the binary if cluster
timing must change.

\subsection{Cluster Membership}

Membership is built from the discovered peers as a stable
\texttt{ClusterConfiguration}. Each peer becomes a voter. Within the running
control state, membership is mutated through the
\texttt{ClusterControlCommand::AddNode} and \texttt{RemoveNode} commands.
Removing a node also drops any writer leases it held.

\subsection{Writer Leases}

A \texttt{WriterLease} grants a node exclusive write authority for an \ipto{}
target. A lease records the \texttt{target}, the \texttt{holder} node, and a
\texttt{LeaseEpoch}. The control state rejects a grant whose holder is not a
known node (\texttt{UnknownLeaseHolder}) and rejects a stale or equal epoch
(\texttt{StaleLeaseEpoch}).

\subsection{Checkpoints}

Two checkpoint records are configured through the control plane: the
\texttt{MetadataOutboxCheckpoint} records the last acknowledged offset for a
(shard, target) pair, and the \texttt{CheckpointManifest} records per-segment
consumption progress for recovery. Both reject stale epochs. Checkpoints are
covered in detail in Section~\ref{sec:recovery}.

\subsection{Environment Variable Summary}

\begin{table}[htbp]
\centering
\begin{tabular}{ll}
\hline
Variable & Used by \\
\hline
\texttt{IPTO\_PG\_HOST} & \ipto{} writer backend \\
\texttt{IPTO\_PG\_PORT} & \ipto{} writer backend \\
\texttt{IPTO\_PG\_USER} & \ipto{} writer backend \\
\texttt{IPTO\_PG\_PASSWORD} & \ipto{} writer backend \\
\texttt{IPTO\_PG\_DATABASE} & \ipto{} writer backend \\
\texttt{VANNAK\_PG\_INTEGRATION} & enables the integration test \\
\texttt{VANNAK\_PG\_PORT} & host port for the test PostgreSQL \\
\texttt{POSTGRES\_USER} & test PostgreSQL container \\
\texttt{POSTGRES\_PASSWORD} & test PostgreSQL container \\
\texttt{POSTGRES\_DB} & test PostgreSQL container \\
\hline
\end{tabular}
\caption{Environment variables recognized across \vannak{}.}
\end{table}

\endinput
